\documentclass[masters,en,listoffigures,listoftables]{NMBU}
% Check the NMBU.cls for all the options and more details

% If you'd like to use chapters instead of sections, add 'book' to the options, as \documentclass[book,...]{NMBU}. Recall that you now have to use \chapter{} and then \section{}, \subsection{}, ... 

% \usepackage[nohyperlinks]{acronym} % If you need a List of Acronyms, consider using this package

% \usepackage[font=small]{caption} % if you prefer smaller captions
% \captionsetup{width=0.9\linewidth} 

% Word of advice: When you are done with your thesis, change the compiler in Overleaf to XeLaTeX! This will make your fonts and title page look much nicer than in pdfLaTeX. Just go to "Menu", in the "Settings", find "Compiler" and change to XeLaTeX.

% The NMBU style package is designed to help you write the cover and back page automatically. Make sure to include the information below to generate your title. You can set the language you are writing in with "en" for english, "bm" for bokmål, or "nn" for nynorsk.

% A table of contents is automatically included. If you have more than one figure, you need a list of figures by adding "listoffigures" to the documentclass parameters. If you have more than one table, you need a list of tables by adding "listoftables" to the documentclass parameters.

% If you would like to have a bibliography with named authors, organised alphabetically, add "namedauthors" to the parameters of the document class.

\addbibresource{references/bib.bib} % here you include the bibliography file. 

% If you are using Norwegian dates in your bibliography, add "norwegiandates" to the document class.

% When you finish your thesis, switch the LaTeX compiler from pdfLaTeX to XelaTeX (on Overleaf, "Menu"->"Compiler"). This is for 2 reasons: 1) to produce your front page with the Arial font; 2) XeLaTeX produces better fonts and (usually) smaller PDFs. Every kilobyte of space we save helps reduce CO2 emissions :)

% Good luck with your thesis! If you get in trouble with LaTeX stuff, write to me: leonardo.rydin@gmail.com

% In your .tex file, you need to include the following information to produce the first page of your thesis

